\section{Discussion}\label{sec:discussion}
\paragraph{What makes the mechanism useful?}
Past utility and present validity are different signals. Frequency can protect heavily used but obsolete records; recency can discard stable facts while keeping recent interference. The state score represents both signals and lets them change on different time scales. In particular, usage records evidence of prior usefulness while its decay limits persistence of that evidence, and consistency can discount a superseded version without assuming that every old fact is obsolete. The component interventions support this interpretation, not a uniquely correct or minimal decomposition.

The rate experiment supplies the strongest evidence because content, weights, and scoring are unchanged within each world. Changing a rate changes the state history that reaches a budget decision. The intervention acts through admission and eviction, which determine available information. The attribution question is whether aggregate replacement volume alone explains performance; the present evidence rejects that explanation in specific mean-matched comparisons, while leaving finer replacement counterfactuals open.

\paragraph{What survives over longer horizons?}
The long-horizon test uses exact 500/1,000/5,000-event prefixes with a context switch every 97 events and renewed observed fact revisions after each switch. At 5,000 events and eight slots, full evolution reaches 75.21\% accuracy versus 43.32\% for fixed decay ($\Delta=31.89$ points, CI $[31.46,32.33]$, Holm $p=0.0090$). Stable retention is 98.43\% versus 45.51\%; stale error is zero for both. At sixteen slots, accuracy is 87.19\% versus 79.37\%, but adaptation is slower by 0.387 target queries (CI $[0.320,0.461]$, adjusted $p=0.0090$). Thus even this favorable long-stream comparison retains a cost.

Frequency's stale-answer rate at sixteen slots increases from 19.09\% at 500 events to 42.39\% at 5,000; full evolution is 0.0042\% at 5,000. The largest observed usage state is 220 without I/C/U flow and approximately 5.29 at strength 1. These are finite-stream measurements, consistent with the contraction of the usage flow. They neither prove infinite-horizon stability nor demonstrate bounded open-world metadata: the stress environment has only 28 fact identities. Full time curves, the collapse diagnostic, and selected-item trajectories are reported in Appendix~\ref{app:stress}.

\paragraph{Stability--plasticity as a policy choice.}
Competing goals here do not establish a universal impossibility theorem: some interventions improve multiple metrics. Practical control would require choosing an acceptable retention loss or stale-error constraint before selecting an operating point, ideally using independent validation. Our grid neither selects deployment strengths automatically nor establishes global Pareto optimality. An equally budgeted sweep of recency--frequency mixtures is an important untested comparator family.

\paragraph{Scope and deployment limitations.}
Structured identities, observed revision signals, and explicit post-answer feedback simplify the synthetic problem. Shared public templates limit independence of content, although schedules are held out. Semantic migration is a controlled replay, not an autonomous tool-using agent. The conversational results use one backend and one same-model judge, too few questions for strong transfer claims, and no competitive full-system study against current agent-memory methods. Storage is constrained in records rather than total bytes; state updates add work. These limitations constrain the paper's claim to a measured mechanism and capability boundary. Deployment additionally needs conflict detection, deletion semantics, provenance, and privacy controls.
